\documentclass[11pt]{article}
\usepackage{preamble}

\newcommand{\IconCheck}{[CHECK]}
\newcommand{\IconMic}{[MIC]}
\newcommand{\IconClipboard}{[CLIPBOARD]}

\newcommand{\phasetag}[1]{%
  {\small\itshape Tag: #1\par}
}

\newcommand{\phaseitems}[1]{%
  \begin{itemize}[leftmargin=1.5em,itemsep=2pt,topsep=2pt]
  #1
  \end{itemize}%
}

\newcommand{\teacheranswerbox}[2]{%
  \begin{callout}{#1}{calloutgreen}
  #2
  \end{callout}
}

\begin{document}

\packetheading{Lesson 4-3 · Period 2 · Teacher Edition}{Multiply \& Divide Rational Expressions + DOK-3 Closure}

\sectionbanner{OBJECTIVES}
{\small
\renewcommand{\arraystretch}{1.25}
\noindent\begin{tabular}{|p{1.8in}|p{4.8in}|}
\hline
\textbf{Math Objective} & Multiply and divide rational expressions by factoring, canceling common factors, and rewriting division as multiplication by the reciprocal; track all domain restrictions through every step. \\ \hline
\textbf{Language Objective} & Justify in writing whether a set is closed under an operation, using sentence frames that cite factoring and cancellation. \\ \hline
\textbf{Essential Understanding} & Rational expressions behave like rational numbers: closed under multiplication and nonzero division, because factoring + cancellation always produces another rational expression. \\ \hline
\end{tabular}
}

\sectionbanner{TOPIC GOALS / ESSENTIAL QUESTION / MATERIALS}
{\small
\renewcommand{\arraystretch}{1.25}
\noindent\begin{tabular}{|p{1.8in}|p{4.8in}|}
\hline
\textbf{Topic Goals} & Multiply and divide rational expressions; justify closure under these operations. \\ \hline
\textbf{Essential Question} & Why are rational expressions closed under multiplication and division, and what must we track to prove it? \\ \hline
\textbf{Materials} & Student packet, pencil. Teacher models Example 3; DOK-3 Practice \#13 is the capstone closure-argument task. \\ \hline
\end{tabular}
}

\begin{callout}{\IconPin\ PRE-FLIGHT — P2 IS THE DOK-3 SPINE DAY}{calloutpink}
Today’s conceptual core is Practice \#13 (Closure Argument). Students must produce a written paragraph with at least ONE worked example showing a product AND one showing a quotient, concluding that the result in each case is itself rational. That argument IS the lesson’s deepest content.\\
Pace: Try-It 3 (4 min) \(\rightarrow\) \#27 (4 min) \(\rightarrow\) Try-It 4 (4 min) \(\rightarrow\) \#29 (4 min) \(\rightarrow\) Try-It 5 (3 min) \(\rightarrow\) \#32 (4 min) \(\rightarrow\) \#13 (12 min). \(\sim\)35 min total Explore.\\
If running tight: cut Try-It 5 or \#32. Never cut \#13.
\end{callout}

\phasetag{bridge from P1}
\frameworkphaseheader{Do Now}{1}{5}{%
\phaseitems{
  \item Hand out at door.
  \item Circulate 3 min.
  \item Cold-call 1 student for factoring and domain restriction.
}%
}{%
\phaseitems{
  \item Simplify by factoring and canceling.
  \item State domain restrictions.
}%
}{%
\phaseitems{
  \item What do you factor first?
  \item Is \(xy\) constant here? What values are excluded from the domain?
}%
}{Solo teacher: circulate silently. No hints.}

\bankitem{Do Now — Simplify this rational expression}{%
What is the simplified form of the rational expression? What is the domain?
\[
\frac{x^3+9x^2-10x}{x^3-9x^2-10x}
\]
}{}

\teacheranswerbox{\IconCheck\ ANSWER KEY}{%
\textbf{Answer key (from source):} \(\dfrac{(x+10)(x-1)}{(x-10)(x+1)}\), \(x \ne -1, 0, 10\).
}

\begin{callout}{\IconTarget\ SENTENCE FRAME}{calloutyellow}
“I factor the numerator as \_\_\_ and the denominator as \_\_\_. I cancel the common factor \_\_\_, leaving \_\_\_. The domain excludes \(x = \_\_\_\) because that value makes the denominator zero.”
\end{callout}

\phasetag{teacher models Example 3}
\frameworkphaseheader{Launch}{2}{12}{%
\phaseitems{
  \item Project Example 3. Think aloud: factor every numerator and denominator before canceling.
  \item “For division I flip the second fraction — the RECIPROCAL — before I do anything else.”
  \item Flag the Multiply/Divide Rules card — print it in student minds.
  \item 30-sec pause: “What domain values did I exclude and why?” (partner check)
  \item Do NOT solve Try-Its or Practice items — release at minute 17.
}%
}{%
\phaseitems{
  \item Watch Example 3 silently.
  \item During pause: state domain restrictions to partner.
  \item Copy Rules card into margin.
  \item Note the division-as-reciprocal step: this comes back in \#13.
}%
}{%
\phaseitems{
  \item What must I do BEFORE canceling in a division problem?
  \item Why can’t I cancel terms (numerator + denominator) — only factors?
  \item If \(x = 2\) makes ANY denominator zero at ANY step, is it excluded? (Yes, always.)
}%
}{Project. Think aloud. Release promptly at minute 17.}

\begin{callout}{\IconMic\ TEACHER SCRIPT}{calloutblue}
1. “Watch me multiply two rational expressions. I factor everything before I cancel.”\\
2. “Step 1: Factor every numerator and denominator completely.”\\
3. “Step 2: For division, rewrite as multiplication by the reciprocal NOW — before factoring or canceling.”\\
4. “Step 3: Cancel only common FACTORS across any numerator with any denominator.”\\
5. “Step 4: State the domain — list every \(x\)-value that makes any denominator zero, including the reciprocal’s new denominator.”\\
6. “Today’s DOK-3 task asks WHY rational expressions are always closed under these operations. You’ll write a paragraph with a worked example.”\\
7. Release to Explore.
\end{callout}

\phasetag{TPS + DOK-3 driver}
\frameworkphaseheader{Explore}{2-3}{33}{%
\phaseitems{
  \item 5 laps across 33 min.
  \item Lap 1 (4 min): Try-It 3 — multiply. Press pairs to factor before canceling.
  \item Lap 2 (4 min): \#27 — multiply with factoring. Watch for premature canceling.
  \item Lap 3 (4 min): Try-It 4 — divide. Check that pairs flip FIRST.
  \item Lap 4 (4 min): \#29 — divide. Confirm domain restriction includes reciprocal’s new denominator.
  \item Lap 5 (12 min): \IconStar\ \#13 Closure Argument. Students write a paragraph with ONE product example AND one quotient example, concluding the result is rational.
  \item Do NOT give \#13 answers. Pairs present argument to each other; one pair presents at Share.
}%
}{%
\phaseitems{
  \item 6 items + \#13 in order: Try-It 3 \(\rightarrow\) \#27 \(\rightarrow\) Try-It 4 \(\rightarrow\) \#29 \(\rightarrow\) Try-It 5 \(\rightarrow\) \#32 \(\rightarrow\) \#13.
  \item For \#13: BEFORE writing, identify what ‘closed’ means and plan your two examples.
  \item Justify with sentence frame.
}%
}{%
\phaseitems{
  \item Did you factor completely before canceling?
  \item \#27: what common factor appears in numerator and denominator?
  \item \#29: what is the reciprocal of the divisor? What new domain restriction does it add?
  \item Try-It 5: both multiply and divide steps — which comes first?
  \item \#13: what makes a set ‘closed’ under an operation? (The result is always in the same set.)
  \item \#13: can you name a product of two rational expressions that is NOT rational? (No — that’s the argument.)
}%
}{Solo teacher: 5 laps. Press CLOSURE argument structure on \#13.}

\bankitem{Try It 3 — Multiply rational expressions}{%
Find the simplified form of each product, and state the domain.

\textbf{a.} \(\dfrac{x^2-16}{9-x}\cdot\dfrac{x^2+x-90}{x^2+14x+40}\)

\textbf{b.} \(\dfrac{x+3}{4x}\cdot\dfrac{3x-18}{6x+18}\cdot\dfrac{x^2}{4x+12}\)
}{}

\teacheranswerbox{\IconCheck\ ANSWER / ANTICIPATED TALK}{%
\textbf{Answer key (from source):} a. \(-(x-4)\) for all \(x\) except \(9\), \(-10\), and \(-4\). b. \(\dfrac{x(x-6)}{32(x+3)}\) for all \(x\) except \(0\) and \(-3\).
}

\bankitem{Practice \#27 — Multiply with factoring}{%
Find the product and the domain.
\[
\frac{(x-y)^2}{x+y}\cdot\frac{3x+3y}{x^2-y^2}
\]
}{}

\teacheranswerbox{\IconCheck\ ANSWER / ANTICIPATED TALK}{%
\textbf{Answer key (from source):} \(\dfrac{3(x-y)}{x+y}\) for all real numbers except \(x=y\) or \(x=-y\).
}

\bankitem{Try It 4 — Divide rational expressions}{%
Find the simplified form of each product and the domain.

\textbf{a.} \(\dfrac{x^3-4x}{6x^2-13x-5}\cdot(2x^3-3x^2-5x)\)

\textbf{b.} \(\dfrac{3x^2+6x}{x^2-49}\cdot(x^2+9x+14)\)
}{}

\teacheranswerbox{\IconCheck\ ANSWER / ANTICIPATED TALK}{%
\textbf{Answer key (from source):} a. \(\dfrac{x^2(x+2)(x-2)(x+1)}{3x+1}\) for all \(x\) except \(\dfrac{5}{2}\) and \(-\dfrac{1}{3}\). b. \(\dfrac{3x(x+2)^2}{x-7}\) for all \(x\) except \(7\) and \(-7\).
}

\bankitem{Practice \#29 — Divide: rewrite as multiplication}{%
Find the product and the domain.
\[
\frac{2x^2-10x}{(x-5)(x^2-1)}\cdot(3x^2+4x+1)
\]
}{}

\teacheranswerbox{\IconCheck\ ANSWER / ANTICIPATED TALK}{%
\textbf{Answer key (from source):} \(\dfrac{2x(3x+1)}{x-1}\) for all \(x\) except \(-1\), \(1\), and \(5\).
}

\bankitem{Try It 5 — Mixed multiply/divide}{%
Find the simplified quotient and the domain of each expression.

\textbf{a.} \(\dfrac{1}{x^2+9x}\div\left(\dfrac{6-x}{3x^2-18x}\right)\)

\textbf{b.} \(\dfrac{2x^2-12x}{x+5}\div\left(\dfrac{x-6}{x+5}\right)\)
}{}

\teacheranswerbox{\IconCheck\ ANSWER / ANTICIPATED TALK}{%
\textbf{Answer key (from source):} a. \(-\dfrac{3}{x+9}\) for all \(x\) except \(0\), \(-9\), and \(6\). b. \(2x\) for all \(x\) except \(6\) and \(-5\).
}

\bankitem{Practice \#32 — Multi-step multiply/divide}{%
Find the quotient and the domain.
\[
\frac{25x^2-4}{x^2-9}\div\frac{5x-2}{x+3}
\]
}{}

\teacheranswerbox{\IconCheck\ ANSWER / ANTICIPATED TALK}{%
\textbf{Answer key (from source):} \(\dfrac{5x+2}{x-3}\) for all \(x\) except \(-3\), \(\dfrac{2}{5}\), and \(3\).
}

\bankitem{Practice \#13 \IconStar\ DOK-3 CLOSURE ARGUMENT}{%
Higher Order Thinking Why are rational expressions closed under multiplication and division?
}{}

\begin{callout}{\IconStar\ DOK-3 CLOSURE ARGUMENT — Practice \#13}{calloutpink}
Students should produce a paragraph with at least ONE worked example showing a product AND one showing a quotient, concluding that the result in each case is itself rational.\\
Key criteria: (1) cites factoring + cancellation by name, (2) shows domain restriction tracking, (3) explicitly concludes the result is a rational expression.
\end{callout}

\teacheranswerbox{\IconCheck\ ANSWER / ANTICIPATED TALK}{%
\textbf{Answer key (from source):} Sample: \(\dfrac{a}{b}\cdot\dfrac{c}{d}=\dfrac{ac}{bd}\), and \(\dfrac{a}{b}\div\dfrac{c}{d}=\dfrac{a}{b}\cdot\dfrac{d}{c}=\dfrac{ad}{bc}\).
}

\phasetag{academic conversation + P3 bridge}
\frameworkphaseheader{Share / Summary}{2}{5}{%
\phaseitems{
  \item Cold-call 1 pair to present their Closure Argument using sentence frame.
  \item Class signals thumbs. Write the closure definition on board.
  \item P3 bridge: “Next period we apply these skills to more complex items — bring your domain-restriction list habit.”
}%
}{%
\phaseitems{
  \item Presenting pair reads their paragraph and cites their two worked examples.
  \item Rest of class signals and writes takeaway in margin.
}%
}{%
\phaseitems{
  \item Summarize: what does it mean for rational expressions to be ‘closed’ under multiplication?
  \item Does division by zero break closure? How do domain restrictions handle that?
}%
}{Cold-call a pair that struggled on \#13 — hold them accountable to the closure-argument structure.}

\begin{sentenceframebox}
\textbf{Sentence frames:}\\
Pair share (Closure Argument): “I rewrite the division as multiplication by the reciprocal, giving \_\_\_. After factoring, I cancel the factor \_\_\_, leaving \_\_\_. The domain excludes \(x = \_\_\_\) because those values make at least one denominator zero at some step.”\\
Whole class: one pair reads the closure argument. Class signals thumbs up/sideways.
\end{sentenceframebox}

\phasetag{summary recap (not graded)}
\frameworkphaseheader{Exit Ticket}{1-2}{3}{%
\phaseitems{
  \item Collect at door.
  \item Scan stem 3 (domain restrictions). Plan P3 Do Now if many students miss the reciprocal’s denominator.
}%
}{%
\phaseitems{
  \item 4 sentence stems, honest.
}%
}{%
\phaseitems{
  \item Did students remember to include reciprocal’s denominator in domain restrictions?
}%
}{Collect. No grade.}

\begin{callout}{\IconClipboard\ EXIT TICKET — Student Form}{calloutyellow}
To divide rational expressions I first \_\_\_.\\
I cancel common factors, not \_\_\_.\\
My domain restriction list must include zeros from \_\_\_.\\
One biggest thing I learned today: \_\_\_\_\_\_\_\_\_\_\_\_\_\_\_\_\_\_\_\_\_\_\_\_\_\_
\end{callout}

\sectionbanner{OPTIONAL REINFORCEMENT (not collected)}
\textit{If you finish Explore early. \#40 is a finance performance task that applies rational expressions to a real-world context.}

\bankitem{Optional \#1  (Savvas Practice, performance-task finance item)}{%
Performance Task The approximate annual interest rate \(r\) of a monthly installment loan is given by the formula:
\[
r=\frac{\left[\dfrac{24(nm-p)}{n}\right]}{p+\dfrac{nm}{12}}
\]

\textbf{Part A} Find the approximate annual interest rate (to the nearest percent) for a four-year signature loan of \(\$20{,}000\) that has monthly payments of \(\$500\).

\textbf{Part B} Find the approximate annual interest rate (to the nearest tenth percent) for a five-year auto loan of \(\$40{,}000\) that has monthly payments of \(\$750\).
}{}

\teacheranswerbox{\IconCheck\ ANSWER / ANTICIPATED TALK}{%
\textbf{Answer key (from source):} Part A \(9\%\). Part B \(4.6\%\).
}

\clearpage

\begin{callout}{\IconClock\ PERIOD A / PERIOD F — 65-MIN CLASS NOTES}{calloutpink}
Both sections should have 65 min for P2. Use the extra 10 min on Explore \#13 (Closure Argument) — the richer the written paragraph, the better the DOK-3 payoff.\\
If finished early: hand out Practice \#40 (finance performance task) as fast-finisher.
\end{callout}

\begin{callout}{\IconPuzzle\ MODIFICATIONS / SUPPORT FOR IEP STUDENTS}{calloutiep}
\textbullet\ Provide the Multiply/Divide Rules card as a sticker/cut-out.\\
\textbullet\ For \#13, provide a paragraph starter: “I will show that the product of two rational expressions is rational. Let \(P(x)/Q(x)\) and \(R(x)/S(x)\) be two rational expressions\ldots”\\
\textbullet\ Accept verbal closure argument in lieu of written paragraph.
\end{callout}

\begin{callout}{\IconGlobe\ SUPPORTS FOR ELL STUDENTS}{calloutell}
\textbullet\ Sentence frames on student page (don’t skip).\\
\textbullet\ Word cards: “reciprocal” = flip the fraction; “closed” = the result stays in the same type; “domain restriction” = values \(x\) cannot equal.\\
\textbullet\ “Closure Argument” = a written explanation showing WHY something is always true.\\
\textbullet\ Pair ELL students with English-dominant partners for TPS.
\end{callout}

\end{document}
